\chapter{Cái giá của được}
\markboth{Cái giá của được}{The Cost of Gaining}

\section{Một người có chức vụ hơn nhưng ngủ ít hơn.}
\begin{truyen}{Một người có chức vụ hơn nhưng ngủ ít hơn.}{Some gains steal your nights.}
\chuhoa{A}{nh từng nghĩ được thăng chức là đích đến. Nhưng khi công việc tràn vào từng bữa cơm và cả giấc ngủ, anh mới hiểu không phải điều gì được thêm cũng làm đời mình đầy lên.}
\textit{(He thought promotion was the destination. But when work filled his meals and even his sleep, he realized that not everything added to life truly enriches it.)}
\end{truyen}
\begin{baihoc}
Có những thứ ta gọi là thành công nhưng thực ra đang ăn mòn phần sống của mình.
\textit{(Some things we call success are quietly eating away at our life.)}
\end{baihoc}
\begin{ghinhoanh}\textbf{Short English to remember:}

Some gains steal your nights.
\end{ghinhoanh}
\begin{tuhocnhanh}
1. Điều em đang theo đuổi lấy đi của em điều gì? \textit{(What is the thing you are chasing taking away from you?)}

2. Cái giá ấy có đáng không? \textit{(Is that price worth it?)}
\end{tuhocnhanh}
\ngancach

\section{Một học sinh được điểm cao nhưng mất niềm vui học.}
\begin{truyen}{Một học sinh được điểm cao nhưng mất niềm vui học.}{A result can cost curiosity.}
\chuhoa{E}{m chỉ học để không rớt hạng. Càng điểm cao, em càng sợ sai. Đến lúc đó, việc học không còn mở ra thế giới nữa mà chỉ còn là cuộc giữ ghế.}
\textit{(The student learned only to avoid falling in rank. The higher the score, the more fearful of mistakes. Learning no longer opened the world, but became a contest to keep a seat.)}
\end{truyen}
\begin{baihoc}
Nếu thành tích làm chết đi sự tò mò, cái được ấy có thể đang quá đắt.
\textit{(If achievement kills curiosity, that gain may be too costly.)}
\end{baihoc}
\begin{ghinhoanh}\textbf{Short English to remember:}

A result can cost curiosity.
\end{ghinhoanh}
\begin{tuhocnhanh}
1. Em đang học vì hiểu hay vì giữ vị trí? \textit{(Are you learning to understand or just to keep your position?)}

2. Niềm vui học của em đang ở đâu? \textit{(Where is your joy in learning right now?)}
\end{tuhocnhanh}
\ngancach

\section{Một người được nhiều lời khen rồi sống để giữ hình ảnh.}
\begin{truyen}{Một người được nhiều lời khen rồi sống để giữ hình ảnh.}{Praise can become a cage.}
\chuhoa{C}{àng được khen là bản lĩnh, anh càng sợ lộ ra phần yếu đuối. Một ngày anh nhận ra mình không còn sống thật, mà đang diễn cho hình ảnh người khác thích.}
\textit{(The more he was praised as strong, the more afraid he became of showing weakness. One day he realized he was no longer living honestly, but performing for an image others liked.)}
\end{truyen}
\begin{baihoc}
Không phải lời khen nào cũng là quà; có lời khen khiến ta quên mất mình thật là ai.
\textit{(Not every compliment is a gift; some make us forget who we really are.)}
\end{baihoc}
\begin{ghinhoanh}\textbf{Short English to remember:}

Praise can become a cage.
\end{ghinhoanh}
\begin{tuhocnhanh}
1. Em có đang sống để giữ hình ảnh nào không? \textit{(Are you living to maintain a certain image?)}

2. Hình ảnh đó có đang làm em mệt không? \textit{(Is that image tiring you?)}
\end{tuhocnhanh}
\ngancach

\section{Một người được tự do hơn nhưng mất phương hướng.}
\begin{truyen}{Một người được tự do hơn nhưng mất phương hướng.}{Freedom without direction can empty you.}
\chuhoa{R}{a khỏi mọi khuôn khổ, cô tưởng mình sẽ hạnh phúc hơn. Nhưng khi không còn biết nên đi đâu, cô mới hiểu tự do đẹp nhất khi đi cùng một hướng sống rõ ràng.}
\textit{(Once free from all structure, she thought she would be happier. But when she no longer knew where to go, she understood that freedom is most beautiful when paired with direction.)}
\end{truyen}
\begin{baihoc}
Có thêm lựa chọn mà không có la bàn cũng là một kiểu lạc lối.
\textit{(Having more choices without a compass is another form of being lost.)}
\end{baihoc}
\begin{ghinhoanh}\textbf{Short English to remember:}

Freedom needs direction.
\end{ghinhoanh}
\begin{tuhocnhanh}
1. Em có đang nhiều lựa chọn nhưng ít định hướng không? \textit{(Do you have many choices but little direction?)}

2. La bàn sống của em là gì? \textit{(What is your compass in life?)}
\end{tuhocnhanh}
\ngancach

\section{Một người được nhiều mối quan hệ nhưng mất sự sâu sắc.}
\begin{truyen}{Một người được nhiều mối quan hệ nhưng mất sự sâu sắc.}{More people, less depth.}
\chuhoa{A}{nh quen rất nhiều người nhưng chẳng còn ai để tâm sự thật. Càng đông bạn, anh càng thấy lẻ loi, và mới hiểu số lượng đôi khi chỉ che giấu sự trống trải.}
\textit{(He knew many people but had no one to truly confide in. The more contacts he had, the lonelier he felt, and he learned that quantity can hide emptiness.)}
\end{truyen}
\begin{baihoc}
Được nhiều chưa chắc là giàu; có khi chỉ là mỏng ra.
\textit{(Having more is not always richness; sometimes it means becoming thinner.)}
\end{baihoc}
\begin{ghinhoanh}\textbf{Short English to remember:}

More people, less depth.
\end{ghinhoanh}
\begin{tuhocnhanh}
1. Em cần thêm người hay thêm chiều sâu? \textit{(Do you need more people or more depth?)}

2. Ai là người em có thể thật lòng với họ? \textit{(Who can you be truly honest with?)}
\end{tuhocnhanh}
\ngancach